\documentclass[11pt]{article}
\usepackage{fontspec}
\setmainfont{Times New Roman}
\setsansfont{Arial}
\setmonofont{Consolas}
\usepackage{amsmath,amssymb,mathtools}
\usepackage{geometry}
\usepackage{microtype}
\newcommand{\mo}{\mathfrak{o}}
\newcommand{\mO}{\mathfrak{O}}
\newcommand{\mT}{\mathfrak{T}}
\newcommand{\mP}{\mathfrak{P}}
\newcommand{\mpideal}{\mathfrak{p}}
\geometry{a4paper,margin=2.35cm}
\setlength{\parindent}{1.25em}
\setlength{\parskip}{0pt}
\makeatletter
\renewcommand\footnotesize{\@setfontsize\footnotesize{7.4}{8.4}\selectfont}
\makeatother
\emergencystretch=100em
\allowdisplaybreaks
\sloppy
\hbadness=10000
\vbadness=10000
\hfuzz=2pt
\vfuzz=10pt
\raggedbottom

\begin{document}

% Унит: Папер31 Сецтион08 Ентры08 дисцриминант басис of екстендед идеал в001. Дирецт Интерславиц Латин транслатион фром цлеан РА34 лине 443 / сегмент P31-С0101, вит сцан-висибле "вхат ис алваыс поссибле", "нацх 1", and the фулл релативе-дисцриминант фоотноте ресторед фром принтед паге 103 / ПДФ паге 22.
\setcounter{footnote}{38}

\noindent Дискриминант $D_{\mT_\mP}$ при том твори базу разширјеного идеала в $\mo_\mpideal$ дискриминантного идеала порјадка $\mT$ взходно к $\mo$. Именно, што јест всегда можно, возмимо елементы $\alpha_1,\ldots,\alpha_n$ из $\mT$ како $\mo_\mpideal$-модулну базу од $\mT_\mP$; тогда $|S(\alpha_i\alpha_k)|$ принадлежи к дискриминантному идеалу, и по пункту 1 всакы детерминант $|S(\tau_i\tau_k)|$ в $\mo_\mpideal$ јест делимы чрез $|S(\alpha_i\alpha_k)|$.\footnote{Из того следује -- ако $\mo$ јест главны порјадок алгебраичного числового польа -- согласје с релативным дискриминантом по Хилберту (кторы, како знано, соглашаје се с определьеньем Hecke). Именно, нехај $\omega_1,\ldots,\omega_m$ буде $\mo$-модулна база главного порјадка релативного польа -- Hilbert особливо избира модулну базу взходно к колцу целых рационалных чисел. Тогда по Хилберту утварјајут се вси детерминанты из всаких $n$ базных елементов и јих конјугованых елементов; "релативны дискриминант" јест идеал в $\mo$, изведены из продуктов всаких двух, једнаких или различных, из тых детерминантов. Тым самым релативны дискриминант в $\mo_\mpideal$ преходи в квадрат детерминанта базы $\alpha_1,\ldots,\alpha_n$ и јеј конјугованых елементов, то јест в $|S(\alpha_i\alpha_k)|$. Зато разширјеньа релативного дискриминанта и дискриминантного идеала (релативного главного порјадка) совпадајут в всих колцах квоциентов $\mo_\mpideal$; по концу пункта 3 оба идеала ставајут идентичне.}

\end{document}

